\section*{Capítulo \ref{ch:g7:areas} --- Áreas y volúmenes}

\begin{solution}{exo:g7:areas:1}
$8 \times 5 = 40$ cm$^2$; \qquad $6.5 \times 4 = 26$ m$^2$.
\end{solution}

\begin{solution}{exo:g7:areas:2}
$A = 10 \times 4 = 40$ cm$^2$. La respuesta no es
$10 \times 6 = 60$ porque el lado de $6$ cm está inclinado: la fórmula
usa la \emph{altura} (la distancia perpendicular), que es $4$ cm.
\end{solution}

\begin{solution}{exo:g7:areas:3}
$\frac{12 \times 7}{2} = 42$ cm$^2$; \quad
$\frac{9 \times 6}{2} = 27$; \quad
$\frac{5.5 \times 2}{2} = 5.5$ m$^2$.
\end{solution}

\begin{solution}{exo:g7:areas:4}
Disco: $\pi \times 5^2 = 25\pi \approx 79$ cm$^2$. Semidisco de \omterm{def:g3:shapes:shapes}{radio}
$10$: $\frac{\pi \times 10^2}{2} = 50\pi \approx 157$ cm$^2$.
\end{solution}

\begin{solution}{exo:g7:areas:5}
\omterm{def:g6:measure:perimeter}{Perímetro}: $2\pi \times 3 = 6\pi \approx 18.8$ cm (una longitud, en cm).
\omterm{def:g6:measure:area}{Área}: $\pi \times 3^2 = 9\pi \approx 28.3$ cm$^2$ (una superficie, en
cm$^2$).
\end{solution}

\begin{solution}{exo:g7:areas:6}
\omterm{def:g7:areas:prism}{Prisma}: $24 \times 10 = 240$ cm$^3$. \omterm{def:g7:areas:prism}{Cilindro}:
$\pi \times 3^2 \times 8 = 72\pi \approx 226$ cm$^3$.
\end{solution}

\begin{solution}{exo:g7:areas:7}
$\frac{8 \times h}{2} = 28$, así que $4h = 28$ y $h = 7$ cm.
\end{solution}

\begin{solution}{exo:g7:areas:8}
Los tres \omterm{def:g2:mult:def}{productos} coinciden porque cada uno calcula la \emph{misma}
magnitud: el \omterm{def:g6:measure:area}{área} del triángulo. Esto da una comprobación práctica
de las construcciones y una manera de calcular una altura a partir de otra
pareja base--altura (se usa en el \cref{exo:g7:areas:12}).
\end{solution}

\begin{solution}{exo:g7:areas:9}
Los dos triángulos tienen \omterm{def:g6:measure:area}{áreas} $\frac{30 \times 20}{2} = 300$ m$^2$
y $\frac{18 \times 20}{2} = 180$ m$^2$: en total, $480$ m$^2$. Fórmula
general sugerida: para lados paralelos $a$ y $b$ a distancia $h$,
\[
A = \frac{(a + b) \times h}{2}
\]
--- la fórmula del trapecio (aquí, $\frac{(30+18)\times 20}{2} = 480$).
\end{solution}

\begin{solution}{exo:g7:areas:10}
\omterm{def:g6:measure:volume}{Volumen} de agua: $\pi \times 6^2 \times 15 = 540\pi \approx 1\,696$
cm$^3$. Base del \omterm{def:g5:solids:def}{ortoedro}: $18 \times 12 = 216$ cm$^2$. Altura alcanzada:
$\frac{540\pi}{216} = 2.5\pi \approx 7.9$ cm.
\end{solution}

\begin{solution}{exo:g7:areas:11}
\omterm{def:g3:shapes:shapes}{Radios} $15$ y $20$ cm. \omterm{def:g6:measure:area}{Áreas}: $225\pi \approx 707$ cm$^2$ y
$400\pi \approx 1257$ cm$^2$. Precio por cada $100$ cm$^2$:
$\frac{9}{7.07} \approx 1.27$ euros y
$\frac{14}{12.57} \approx 1.11$ euros. La pizza grande sale mejor de
precio --- las \omterm{def:g6:measure:area}{áreas} crecen con el \emph{cuadrado} del \omterm{def:g4:geometry:circle}{diámetro}.
\end{solution}

\begin{solution}{exo:g7:areas:12}
Con los catetos: $A = \frac{6 \times 8}{2} = 24$ cm$^2$. Con la
hipotenusa como base: $24 = \frac{10 \times h}{2} = 5h$, así que
$h = 4.8$ cm.
\end{solution}

\begin{solution}{pb:g7:areas:1}
\textbf{1.} Cuadrado: $8 \times 8 = 64$. Rectángulo:
$5 \times 13 = 65$. Las mismas cuatro piezas parecen cubrir $64$
unidades cuadradas en una figura y $65$ en la otra: se ha creado una unidad
de \omterm{def:g6:measure:area}{área} de la nada --- supuestamente.

\textbf{2.} Cada triángulo: $\frac{8 \times 3}{2} = 12$. Cada
trapecio: $\frac{(5 + 3) \times 5}{2} = 20$.

\textbf{3.} $12 + 12 + 20 + 20 = 64$: las piezas suman $64$,
así que pueden llenar exactamente el cuadrado --- y \emph{no} pueden llenar
el rectángulo de $65$ unidades. El dibujo del rectángulo tiene que contener un
hueco.

\textbf{4.} La misma inclinación significaría que $3$ en $8$ es
proporcional a $2$ en $5$; \omterm{def:g2:mult:def}{productos} cruzados:
$3 \times 5 = 15$ y $2 \times 8 = 16$. No son iguales: la
hipotenusa y el borde inclinado \emph{no} se alinean. La
«diagonal» del rectángulo es mentira --- dos pendientes ligeramente
distintas que se juntan en un ángulo muy abierto.

\textbf{5.} A lo largo de la falsa diagonal, las cuatro piezas dejan un
hueco larguísimo y finísimo --- una astilla con forma de
\omterm{def:g7:central:parallelogram}{paralelogramo} muy achatado --- cuya \omterm{def:g6:measure:area}{área} es exactamente la unidad
cuadrada que falta, $65 - 64 = 1$. Estirada a lo largo de una diagonal de unas
$13$ unidades, su anchura media es de unas
$1 \div 13 \approx 0.08$ unidades: en un dibujo real, más fina que
el trazo de lápiz que la esconde.

\textbf{6.} Cada triángulo tiene base $6$ cm y altura $4$ cm ---
la distancia entre las \omterm{def:g6:lines:perp}{paralelas}, esté donde esté el \omterm{def:g5:solids:def}{vértice} ---,
así que cada \omterm{def:g6:measure:area}{área} es $\frac{6 \times 4}{2} = 12$ cm$^2$
(\cref{thm:g7:areas:triangle}). Deslizar el \omterm{def:g5:solids:def}{vértice} a lo largo de la
paralela cambia la forma, nunca la base ni la altura:
\omterm{def:g6:measure:area}{áreas} iguales para siempre.

\textbf{7.} La diagonal corta el trapecio en un triángulo
de base $B$ y otro de base $b$, los dos de altura $h$ (la
distancia entre los lados paralelos):
\[
A = \frac{B \times h}{2} + \frac{b \times h}{2}
= \frac{(B + b) \times h}{2} .
\]

\textbf{8.} $A = \frac{(9 + 5) \times 4}{2} = \frac{56}{2} =
28$ cm$^2$.

\textbf{9.} La media de los lados paralelos es
$\frac{B + b}{2}$, así que (media) $\times h =
\frac{(B+b) \times h}{2}$: la misma fórmula, factorizada
de otro modo. Comprobación: media $= \frac{9+5}{2} = 7$, y
$7 \times 4 = 28$ cm$^2$.

\textbf{10.} La pila torcida contiene exactamente las mismas cartas
--- las mismas capas, cada una de la misma \omterm{def:g6:measure:area}{área} ---, así que el mismo
\omterm{def:g6:measure:volume}{volumen}. En $V = B \times h$, la altura hay que medirla
\emph{perpendicularmente} a la base (la altura vertical de la
pila), no a lo largo de la inclinación: exactamente igual que el \omterm{def:g6:measure:area}{área} del
\omterm{def:g7:central:parallelogram}{paralelogramo} quiere la altura perpendicular y no el lado inclinado
(\cref{exo:g7:areas:2}), una dimensión más arriba.

\textbf{11.} Corona $= \pi \times 5^2 - \pi \times 3^2 = 25\pi -
9\pi = 16\pi \approx 50$ m$^2$ (\cref{thm:g7:areas:disk}).

\textbf{12.} Un disco de \omterm{def:g6:measure:area}{área} $16\pi$ tiene \omterm{def:g3:shapes:shapes}{radio} $4$, ya que
$\pi \times 4^2 = 16\pi$. Los \omterm{def:g3:shapes:shapes}{radios} $3$, $4$ y $5$: el trío más
famoso de las matemáticas, protagonista del
\cref{ch:g8:pythagoras}.

\textbf{13.} Una pizza de $40$ cm: \omterm{def:g3:shapes:shapes}{radio} $20$ y
\omterm{def:g6:measure:area}{área} $\pi \times 400 \approx 1\,257$ cm$^2$. Dos pizzas de $28$ cm:
\omterm{def:g3:shapes:shapes}{radio} $14$ y, juntas,
$2 \times \pi \times 196 = 392\pi \approx 1\,232$ cm$^2$. La
pizza grande sola es \emph{más} pizza que las dos medianas.

\textbf{14.} Multiplicar el \omterm{def:g3:shapes:shapes}{radio} por $1.4$ multiplica el \omterm{def:g6:measure:area}{área} por
$1.4^2 = 1.96$ --- casi el doble. El factor exacto es el
número cuyo cuadrado es $2$, unos $1.414$: el número de la diagonal
del \cref{ex:g8:pythagoras:diagonal}. (Regla práctica: una \omterm{def:g3:division:half}{mitad} más
de anchura es casi el doble de pizza.)

\textbf{15.} $13 \times 13 = 169$, pero $8 \times 21 = 168$:
esta vez \emph{desaparece} una unidad cuadrada --- las piezas
se solapan en una astilla fina en vez de dejar un hueco. Con
términos consecutivos $3, 5, 8, 13, 21$ de la sucesión de
Hemachandra--Fibonacci, los \omterm{def:g2:mult:def}{productos} van alternando entre uno
más y uno menos que el cuadrado ($5 \times 13 = 65 =
8^2 + 1$, y después $8 \times 21 = 168 = 13^2 - 1$): el mago
«gana» y «pierde» alternativamente una unidad. Las \omterm{def:g6:measure:area}{áreas} no mienten nunca:
unas piezas que suman $64$ cubren $64$, las deslices como las deslices --- toda
unidad que falte o sobre se esconde en una astilla, esperando a que una
comprobación de pendientes la delate.
\end{solution}
